\section{Six domains, one deployable}
\label{sec:ch02-six-domains}

The \code{Product} that chapter~\ref{ch:why-one-graph-is-never-enough} sketched
had eight fields on it and four owners: merchandising writes the title,
pricing sets what it costs today, the warehouse counts what is on the shelf, and
customers supply the opinions. That was presented as the problem. Here it is as
a directory listing:

\begin{minted}{text}
src/Mosaic.Api/
  Catalog/    Pricing/    Inventory/
  Reviews/    Accounts/   Ordering/
  Infrastructure/
\end{minted}

\index{Mosaic!domain folders}%
Each folder has the same three subfolders: \code{Model} for the records,
\code{Data} for the store and the service, \code{Types} for the GraphQL surface.
The rule that makes this worth doing is not the folder structure, though. It is
this: \emph{a field lives in the folder of the domain that owns it, even when it
hangs off a type another domain declared}.

\code{Product} is a Catalog record. \code{Product.price} is not a property on it.
It is a resolver in Pricing's folder:

\begin{minted}{csharp}
[ObjectType<Product>]
public static partial class ProductPricingNode
{
    public static async Task<Money> GetPriceAsync(
        [Parent] Product product,
        PricingService pricing,
        CancellationToken cancellationToken)
    {
        var price = await pricing.GetPriceByProductIdAsync(product.Id, cancellationToken);

        return price?.Amount
            ?? throw new InvalidOperationException(
                $"No price is seeded for product {product.Id}.");
    }
}
\end{minted}

Inventory does the same for the stock count, and Reviews for the review list and
the average rating. Four classes in four folders contribute to one
\code{Product}, and Ordering does the same trick in the other direction,
attaching \code{product} to its own \code{OrderLine} type by reaching into
Catalog.

\index{HotChocolate!type extensions}%
This works because \code{[ObjectType<T>]} is not exclusive. Several classes in
one assembly can extend the same type, and the source generator merges them.
I want to flag that I did not take this on faith: the documented examples of the
pattern put each extension in a separate assembly, which is what a federated
service looks like, and I could not find it stated anywhere that a single
assembly may do it too. So I wrote a throwaway project with two
\code{[ObjectType<Product>]} classes in it before committing the layout to a
book. It compiles, and both fields appear.

The generated SDL is one type assembled from four folders:

\begin{graphqlsdl}
type Product {
  id: ID!
  availableQuantity: Int!
  price: Money!
  reviews: [Review!]!
  averageRating: Float
  sku: String!
  title: String!
  description: String
  category: ProductCategory!
}
\end{graphqlsdl}

\index{HotChocolate!field ordering}%
The field order is worth a glance, because it is not the order anybody wrote.
Fields declared by extension classes come first, then the record's own inferred
properties. If you are diffing exported schemas in code review, and
section~\ref{sec:ch02-dev-loop} argues you should be, expect a new resolver to
appear in the middle of a type rather than at the end.

\subsection*{Why this layout and not the obvious one}

The obvious layout is \code{Models/}, \code{Services/} and \code{Types/} at the
project root, with everything of a kind together. It is what most .NET samples
do and there is nothing wrong with it for a service that will stay one service.

Mosaic will not stay one service, and that changes the arithmetic.
Chapter~\ref{ch:the-first-cut-extracting-catalog} lifts Catalog out into its own
subgraph. Under the layout above that is a \code{git mv} of one folder, a new
project, and a four-line \code{Product} stub in the remaining monolith carrying
nothing but the key. Not one resolver body changes, because Pricing's price
resolver was never in Catalog's file to begin with. Under the type-per-folder
layout, the same extraction means opening three shared folders and deciding, file
by file, which half of each one is leaving. I have done the second kind of split
and would rather not describe it in a chapter I want people to enjoy.

\index{Mosaic!domain dependencies}%
The dependencies between domains also all point one way. Pricing, Inventory,
Reviews and Ordering reference Catalog's \code{Product}. Reviews and Ordering
reference Accounts' \code{Customer}. Catalog references nobody. That acyclic
shape is why Catalog is the domain that can be extracted first, and
chapter~\ref{ch:the-first-cut-extracting-catalog} collects
on it. If you are laying out your own service with a future split in mind, the
question to ask of each dependency is not whether it is tidy but whether it
points at something you could deploy separately tomorrow.

\begin{figure}[htbp]
  \centering
  % Mosaic at the ch02 tag: six domains in one process, and what is not there
% yet. Referenced from chapter 02. Bare tikzpicture: the figure environment,
% caption and label live at the call site. Uses only arrows.meta and positioning.
\begin{tikzpicture}[
    font=\small,
    dom/.style={draw, rounded corners=2pt, minimum width=25mm,
                minimum height=9mm, align=center},
    store/.style={draw, rounded corners=2pt, minimum width=62mm,
                  minimum height=9mm, align=center},
    entry/.style={draw, thick, rounded corners=2pt, minimum width=34mm,
                  minimum height=9mm, align=center},
    consumer/.style={draw, dotted, rounded corners=2pt, minimum width=24mm,
                     minimum height=8mm, align=center},
    later/.style={draw=gray, dashed, text=gray, rounded corners=2pt,
                  minimum width=27mm, minimum height=11mm, align=center},
    flow/.style={-{Stealth[length=2mm]}, thick},
    note/.style={font=\scriptsize\itshape, align=center},
    laternote/.style={font=\scriptsize\itshape, align=center, text=gray}
  ]

  \node[consumer] (client) at (3.0,5.1) {a client};
  \node[entry]    (entry)  at (3.0,3.5) {\texttt{/graphql}};

  \draw[flow] (client) -- (entry);
  \draw[flow] (entry.south) -- (3.0,2.2);

  \node[dom] (cat) at (-3.1,1.35) {Catalog};
  \node[dom] (pri) at ( 0,   1.35) {Pricing};
  \node[dom] (inv) at ( 3.1, 1.35) {Inventory};
  \node[dom] (rev) at (-3.1, 0.25) {Reviews};
  \node[dom] (acc) at ( 0,   0.25) {Accounts};
  \node[dom] (ord) at ( 3.1, 0.25) {Ordering};

  \node[store] (data) at (0,-1.15) {seeded in-memory data};
  \node[note, anchor=north] at (0,-1.8)
    {chapter 4 replaces this with a real database};

  % The one process everything ships in.
  \draw[thick] (-4.75,2.2) rectangle (4.75,-2.35);
  \node[anchor=south west, font=\small\itshape] at (-4.75,2.2)
    {Mosaic v1: one process, one deployable};

  % Deliberately absent at this tag.
  \node[laternote, anchor=north] at (0,-2.6)
    {dashed grey: not in v1, and deliberately so};

  \node[later] (db)     at (-3.1,-3.7) {a real\\database};
  \node[later] (router) at ( 0,  -3.7) {a router\\in front};
  \node[later] (split)  at ( 3.1,-3.7) {separate\\services};

  \node[laternote, anchor=north] at (-3.1,-4.45) {chapter 4};
  \node[laternote, anchor=north] at ( 0,  -4.45) {chapter 10};
  \node[laternote, anchor=north] at ( 3.1,-4.45) {chapters 8 and 12};

\end{tikzpicture}

  \caption{Mosaic at the end of this chapter. Six domains, one process, one
  endpoint, and a seeded list where a database will go. The dashed additions are
  the ones the book spends Parts~I to~III putting in, and the seams they follow
  are the folder boundaries that already exist.}
  \label{fig:ch02-mosaic-v1}
\end{figure}

\subsection*{Registration}

Program.cs ends up reading as a manifest of what Mosaic is made of:

\begin{minted}{csharp}
builder.Services
    .AddCatalogDomain()
    .AddPricingDomain()
    .AddInventoryDomain()
    .AddReviewsDomain()
    .AddAccountsDomain()
    .AddOrderingDomain();

builder.AddGraphQL().AddMosaic();
\end{minted}

Six lines of dependency injection, one line of GraphQL. The asymmetry is the
point of the module attribute from section~\ref{sec:ch02-smallest-service}:
because the generator works across the whole assembly, adding a domain folder
does not mean adding a GraphQL registration. It means adding a service
registration and nothing else.

That list is also the file
chapter~\ref{ch:the-first-cut-extracting-catalog} cuts in half.

Mosaic is now a service rather than an argument. What it is not yet is
something you can work on comfortably, and that is the next thing to fix.
\index{Mosaic|)}
